
\section*{Quiche Lorraine - der Klassiker}
\ihead{}\chead{Personen: 4}\ohead{}
\ifoot{Vorbereitungszeit: 25 Min}\cfoot{}\ofoot{Kochzeit: 35 Min}
{\Large Zutaten}
\begin{multicols}{2}
\textit{Mürbeteig:}
\begin{itemize}
    \item \SI{250}{g} Mehl
    \item \SI{160}{g} eiskalte Butter
    \item \num{1} Ei
    \item \num{1} Prise Salz
    \item etwas Mehl für die Arbeitsfläche
\end{itemize}
\textit{Belag:}
\begin{itemize}
    \item \SI{100}{g} Speckwürfel
    \item ca. \SI{15}{g} Butter zum Braten
    \item \num{1} Zwiebel, gewürfelt
    \item \num{1} Stange Porree
    \item \num{4} große Eier
    \item \SI{250}{ml} Sahne
    \item \num{1} Prise Pfeffer
\end{itemize}
\end{multicols}
\noindent
{\Large Zubereitung}\\
\\
\textit{Vorbereitung:} Für eine Quicheform mit \SIrange{28}{30}{cm} Durchmesser. Ofen auf \SI{180}{\celsius} Ober-/Unterhitze vorheizen.\\
Mehl auf die Arbeitsfläche häufen und eine Mulde in die Mitte drücken.\\
Das Ei, die eiskalte Butter in Flocken und eine Prise Salz hineingeben.\\
Alle Zutaten schnell zu einem Teig verkneten.\\
Den Teig in Frischhaltefolie wickeln und für \SI{30}{min} kaltstellen.\\
Die Speckwürfel in etwas Butter anbraten.\\
Den Porree putzen, in feine Ringe schneiden und zusammen mit den Zwiebelwürfeln zum Speck geben.\\
Den Mürbeteig auf einer bemehlten Arbeitsfläche ausrollen und in die Quicheform geben.\\
Den Teig in der Form festdrücken.\\
Die Eier und die Sahne mit Pfeffer verrühren und Porree, Zwiebeln und Speck unterheben.\\
Die Füllung auf dem Teig verteilen.\\
Die Quiche im vorgeheizten Ofen bei \SI{180}{\celsius} Ober-/Unterhitze für ca. \SI{35}{min} backen.
